\documentclass[12pt,a4paper]{report}

\usepackage[utf8]{inputenc}
\usepackage[T1]{fontenc}
\usepackage[french]{babel}
\usepackage{amsmath, amssymb}
\usepackage[ruled,vlined]{algorithm2e}
\usepackage{graphicx}
\usepackage{booktabs}
\usepackage{listings}
\usepackage{xcolor}
\usepackage{hyperref}

\lstset{
	language=Python,
	basicstyle=\ttfamily\small,
	keywordstyle=\color{blue},
	commentstyle=\color{gray},
	stringstyle=\color{red},
	frame=single,
	breaklines=true,
	literate=
		{é}{{\'e}}1 {è}{{\`e}}1 {ê}{{\^e}}1
		{à}{{\`a}}1 {â}{{\^a}}1
		{ù}{{\`u}}1 {û}{{\^u}}1
		{î}{{\^i}}1 {ï}{{\"i}}1
		{ô}{{\^o}}1
		{ç}{{\c{c}}}1
}

\title{Documentation -- Slime Mold Algorithm appliqué au TSP}
\author{Quentin Beauquier}
\date{\today}

\begin{document}

\maketitle
\tableofcontents
\newpage

% -----------------------------------------------
\chapter*{Nomenclature}
\addcontentsline{toc}{chapter}{Nomenclature}

\section*{Symboles mathématiques}

\begin{table}[h]
	\centering
	\begin{tabular}{@{}cl@{}}
		\toprule
		\textbf{Symbole} & \textbf{Définition} \\
		\midrule
		$X_i$     & Tournée (liste ordonnée de sommets) de l'agent $i$ \\
		$X_b$     & Tournée de la meilleure solution trouvée \\
		$X_A, X_B$ & Tournées de deux agents sélectionnés par tirage pondéré \\
		$W(X_i)$  & Poids attribué à l'agent $X_i$, fonction de son rang de fitness \\
		$S(X_i)$  & Valeur de fitness (coût total de la tournée) de l'agent $X_i$ \\
		$bF$      & Meilleure fitness de l'itération courante (\textit{best fitness}) \\
		$wF$      & Pire fitness de l'itération courante (\textit{worst fitness}) \\
		$p$       & Probabilité de croisement : $p = \tanh(|S(X_i) - bF|)$ \\
		$r$       & Nombre aléatoire tiré uniformément dans $[0, 1]$ \\
		$n$       & Nombre total d'agents dans la population \\
		$t$       & Indice de l'itération courante \\
		$T$       & Nombre maximal d'itérations \\
		\bottomrule
	\end{tabular}
\end{table}

\section*{Termes techniques}

\begin{table}[h]
	\centering
	\begin{tabular}{@{}lp{10cm}@{}}
		\toprule
		\textbf{Terme} & \textbf{Définition} \\
		\midrule
		Métaheuristique & Méthode d'optimisation générique, souvent stochastique, applicable à une large classe de problèmes sans garantie d'optimalité \\
		Myxomycète & Organisme unicellulaire capable de former des réseaux de transport efficaces ; \textit{Physarum polycephalum} en est l'espèce de référence \\
		Fitness & Fonction objectif évaluant la qualité d'une solution ; dans le cas du TSP, il s'agit du coût total de la tournée (somme des distances euclidiennes) \\
		Agent & Entité de la population représentant une solution candidate (une tournée complète), analogue à un filament du myxomycète \\
		Tournée & Permutation des sommets du graphe définissant un cycle hamiltonien \\
		Exploitation & Phase où l'algorithme affine les solutions proches de la meilleure solution connue \\
		Exploration & Phase où l'algorithme diversifie la recherche dans des régions inexplorées de l'espace des solutions \\
		OX Crossover & Opérateur de croisement ordonné (\textit{Order Crossover}) préservant l'ordre relatif des éléments entre deux solutions parentes \\
		Random Swap & Opérateur de mutation échangeant aléatoirement deux sommets dans une tournée \\
		2-opt & Heuristique d'amélioration locale consistant à supprimer deux arêtes d'une tournée et à les reconnecter de manière à réduire le coût total \\
		TSP & Problème du voyageur de commerce (\textit{Travelling Salesman Problem}) : déterminer le cycle hamiltonien de coût minimal dans un graphe complet \\
		Gap & Écart relatif entre le coût d'une solution obtenue et le coût de la solution optimale connue \\
		Convergence & État atteint lorsque l'algorithme ne parvient plus à améliorer significativement la solution au fil des itérations \\
		Arrêt anticipé & Mécanisme interrompant l'exécution lorsque ni la phase de 2-opt ni la phase d'exploration ne produisent d'amélioration \\
		\bottomrule
	\end{tabular}
\end{table}

\newpage

% -----------------------------------------------
\chapter{Définition}

Le \textit{Slime Mold Algorithm} (SMA) est une métaheuristique d'optimisation
inspirée du comportement du myxomycète \textit{Physarum polycephalum}.
Proposé par Li et al. (2020), il modélise la capacité de cet organisme
à former des réseaux de transport efficaces en explorant son environnement
à la recherche de nourriture.

\section{Principe de base}
Le \textit{SMA} simule le comportement du myxomycète.
Celui-ci se déploie de manière stochastique afin d'explorer son environnement
et de localiser des sources nutritives, en établissant des canaux permettant
leur acheminement.
Lorsqu'il rencontre une source de nourriture, le myxomycète renforce le chemin le
plus court afin d'optimiser sa consommation d'énergie et supprime les
canaux inutiles.

\section{Adaptation au TSP}
L'algorithme original de Li et al. opère dans un espace continu. Notre
implémentation l'adapte au problème du voyageur de commerce, qui relève
de l'optimisation combinatoire. Les mises à jour de position continues
sont remplacées par des opérateurs discrets sur des permutations :
\begin{itemize}
	\item Le \textbf{croisement ordonné} (OX) se substitue à la mise à jour
	      vectorielle de la phase d'exploitation ;
	\item La \textbf{mutation par échange aléatoire} (\textit{random swap})
	      remplace la contraction de la phase d'exploration ;
	\item L'optimisation locale \textbf{2-opt} est intégrée de manière cyclique
	      afin d'affiner les tournées obtenues.
\end{itemize}

% -----------------------------------------------
\chapter{Architecture}

L'implémentation se compose de quatre modules principaux.

\section{Graph}
La classe \texttt{Graph} encapsule le problème TSP. Elle stocke la liste
des coordonnées des villes et fournit les méthodes suivantes :
\begin{itemize}
	\item \texttt{distance(i, j)} : distance euclidienne entre les villes $i$ et $j$ ;
	\item \texttt{turn\_cost(turn)} : coût total d'une tournée (somme des distances
	      consécutives, y compris le retour à la ville de départ) ;
	\item \texttt{random(nb\_cities)} : génération d'une instance aléatoire ;
	\item \texttt{plot(turn, ...)} : visualisation de la tournée et de la courbe
	      de convergence.
\end{itemize}

\section{Agent}
La classe \texttt{Agent} représente une solution candidate. Chaque agent
possède trois attributs :
\begin{itemize}
	\item \texttt{turn} : la tournée courante (permutation des indices de villes) ;
	\item \texttt{cost} : le coût de cette tournée, calculé à l'initialisation
	      via \texttt{graph.turn\_cost} ;
	\item \texttt{weight} : le poids utilisé lors de la sélection pondérée
	      (initialisé à $1{,}0$).
\end{itemize}
La méthode de classe \texttt{Agent.random(graph)} crée un agent dont la
tournée est une permutation aléatoire de l'ensemble des villes.

\section{Opérateurs}
Trois opérateurs sont définis dans le module \texttt{operators} :

\subsection{Random Swap (mutation)}
L'opérateur \texttt{random\_swap} sélectionne deux positions aléatoires
dans la tournée et permute les villes correspondantes. Il s'agit de
l'opérateur d'exploration, appliqué lorsqu'un agent est jugé trop
éloigné de la meilleure solution.

\subsection{Order Crossover -- OX (croisement)}
L'opérateur \texttt{ox\_crossover} effectue un croisement ordonné entre
deux parents. Il sélectionne un segment aléatoire $[a, b)$ du premier
parent, le copie dans l'enfant, puis complète les positions restantes
avec les villes du second parent, en préservant leur ordre relatif
d'apparition.

Dans notre implémentation, le croisement est appliqué en cascade :
\begin{equation}
	X' = \text{OX}\bigl(X_b,\; \text{OX}(X_A,\; X_B)\bigr)
\end{equation}
Le croisement intermédiaire $\text{OX}(X_A, X_B)$ recombine deux agents
sélectionnés par tirage pondéré, puis le résultat est croisé avec la
meilleure solution $X_b$, orientant ainsi la recherche vers les régions
les plus prometteuses.

\subsection{2-opt (amélioration locale)}
L'opérateur \texttt{two\_opt} parcourt toutes les paires d'arêtes
$(i, i{+}1)$ et $(j, j{+}1)$ de la tournée. Si la suppression de ces
deux arêtes et leur reconnexion croisée diminue le coût total, le
segment $[i{+}1, j]$ est inversé. L'opérateur effectue au plus
\texttt{max\_passes} passes complètes et s'arrête prématurément si
aucune amélioration n'est trouvée lors d'une passe.

\section{Population}
La classe \texttt{Population} orchestre l'ensemble de l'algorithme. Elle
gère la collection d'agents et expose deux méthodes principales :
\texttt{step} (une itération) et \texttt{run} (boucle complète avec
arrêt anticipé).

% -----------------------------------------------
\chapter{Méthode}

\section{Calcul des poids}
À chaque itération, les agents sont triés par coût croissant. Le poids
$W(X_i)$ de chaque agent est calculé en fonction de son rang :

\begin{equation}
	W(X_i) =
	\begin{cases}
		1 + r \cdot \log\!\left(\dfrac{bF - S(X_i)}{bF - wF} + 1\right)
		& \text{si } i < \lfloor n/2 \rfloor \\[6pt]
		1 - r \cdot \log\!\left(\dfrac{bF - S(X_i)}{bF - wF} + 1\right)
		& \text{si } i \geq \lfloor n/2 \rfloor
	\end{cases}
\end{equation}

où $r \sim \mathcal{U}(0, 1)$. Les agents de la moitié supérieure
(meilleures solutions) reçoivent un poids supérieur à $1$, renforçant
leur probabilité d'être sélectionnés lors du croisement. Les agents de
la moitié inférieure reçoivent un poids inférieur à $1$. Si $bF = wF$
(population homogène), les poids demeurent inchangés.

\section{Mise à jour des positions}
Pour chaque agent, la probabilité de croisement est calculée :
\begin{equation}
	p = \tanh\!\bigl(|S(X_i) - bF|\bigr)
\end{equation}

Plus un agent est éloigné de la meilleure solution, plus $p$ est élevé,
et plus la probabilité de croisement augmente. La mise à jour s'effectue
selon :

\begin{equation}
	X_i(t+1) =
	\begin{cases}
		\text{OX}\bigl(X_b,\; \text{OX}(X_A,\; X_B)\bigr)
		& \text{si } r < p \quad \text{(exploitation)} \\[4pt]
		\text{random\_swap}(X_i)
		& \text{sinon} \quad \text{(exploration)}
	\end{cases}
\end{equation}

où $X_A$ et $X_B$ sont sélectionnés par tirage pondéré (sans remplacement
entre eux) proportionnellement à leur poids $W$.

La nouvelle tournée n'est acceptée que si son coût est strictement
inférieur au coût courant de l'agent (\textbf{acceptation gloutonne}).

\section{Phases cycliques}
L'exécution est structurée en \textbf{cycles} de longueur
\texttt{phase\_length} (par défaut $15$) itérations. Chaque cycle se
décompose en deux phases :

\begin{enumerate}
	\item \textbf{Phase d'exploration}
	      (premières $\lceil(1 - \texttt{opt\_ratio}) \times \texttt{phase\_length}\rceil$
	      itérations) : seuls les opérateurs de croisement et de mutation sont
	      appliqués.
	\item \textbf{Phase d'optimisation locale}
	      (dernières $\lfloor\texttt{opt\_ratio} \times \texttt{phase\_length}\rfloor$
	      itérations, par défaut $25\,\%$) : l'opérateur 2-opt est
	      appliqué en complément à chaque agent, avec au plus
	      \texttt{max\_opt\_passes} passes (par défaut $5$).
\end{enumerate}

Cette alternance permet de diversifier la recherche avant de l'affiner
localement, évitant ainsi une convergence prématurée.

\section{Arrêt anticipé}
Le mécanisme d'arrêt anticipé surveille la stagnation de l'algorithme :
\begin{enumerate}
	\item À l'entrée de chaque phase 2-opt, le meilleur coût courant est
	      enregistré (\textit{snapshot}).
	\item Si, à l'issue de la phase 2-opt, le meilleur coût n'a pas diminué
	      (phase jugée \textit{stale}), un indicateur est levé.
	\item Si, lors de la phase d'exploration suivante, le meilleur coût
	      demeure inchangé, l'algorithme conclut que ni l'exploitation ni
	      l'exploration ne produisent d'amélioration et met fin à l'exécution.
\end{enumerate}

% -----------------------------------------------
\chapter{Analogie biologique}
Le myxomycète explore son environnement en déployant ses filaments vers
les sources de nourriture. Les filaments menant aux meilleures sources
se renforcent, les autres s'atrophient.

Dans l'algorithme, chaque agent représente un filament. Les agents proches
de bonnes solutions voient leur poids $W$ augmenter (phase d'exploitation),
tandis que les agents éloignés explorent davantage (phase d'exploration).
La probabilité $p = \tanh(|S(X_i) - bF|)$ traduit ce mécanisme : plus un
filament est éloigné d'une source nutritive, plus il est susceptible d'être
reconfiguré par croisement avec des filaments mieux positionnés.

\chapter{Cas d'application connus}

\begin{itemize}
	\item Optimisation de fonctions continues (benchmark CEC)
	\item Sélection de caractéristiques en apprentissage automatique
	\item Problèmes d'ingénierie (conception mécanique, réseaux électriques)
	\item Problème du voyageur de commerce (TSP)
	\item Ordonnancement de tâches
\end{itemize}

Référence principale : Li, S., Chen, H., Wang, M., Heidari, A. A., \& Mirjalili, S. (2020).
\textit{Slime mould algorithm: A new method for stochastic optimization}.
Future Generation Computer Systems, 111, 300--323.

% -----------------------------------------------
\chapter{Hyperparamètres}

\begin{table}[h]
	\centering
	\begin{tabular}{@{}llcl@{}}
		\toprule
		\textbf{Paramètre} & \textbf{Description} & \textbf{Défaut} & \textbf{Unité} \\
		\midrule
		\texttt{nb\_agents}      & Taille de la population                & 40  & agents \\
		\texttt{T}               & Nombre maximal d'itérations            & 200 & itérations \\
		\texttt{phase\_length}   & Longueur d'un cycle exploration/2-opt  & 15  & itérations \\
		\texttt{opt\_ratio}      & Proportion du cycle dédiée au 2-opt    & 0,25 & -- \\
		\texttt{max\_opt\_passes}& Passes maximales de 2-opt par itération & 5   & passes \\
		\bottomrule
	\end{tabular}
	\caption{Hyperparamètres de l'algorithme et valeurs par défaut}
\end{table}

% -----------------------------------------------
\chapter{Exemples de code}

\section{Agent}
\begin{lstlisting}
class Agent:
    def __init__(self, graph, turn):
        self.graph = graph
        self.turn = turn
        self.weight = 1.0
        self.cost = self.graph.turn_cost(turn)

    @classmethod
    def random(cls, graph):
        turn = list(range(len(graph.cities)))
        random.shuffle(turn)
        return cls(graph, turn)
\end{lstlisting}

\section{Calcul des poids}
\begin{lstlisting}
def update_weights(self):
    n = len(self.agents)
    self.agents.sort(key=lambda a: a.cost)
    bF = self.best.cost
    wF = self.worst.cost
    if bF == wF:
        return
    for i in range(0, n):
        s = self.agents[i].cost
        calc = np.log((bF - s) / (bF - wF) + 1)
        if i < n / 2:
            self.agents[i].weight = 1 + random.uniform(0, 1) * calc
        else:
            self.agents[i].weight = 1 - random.uniform(0, 1) * calc
\end{lstlisting}

\section{Itération principale}
\begin{lstlisting}
def step(self, t=0, phase_length=20, opt_ratio=0.25, max_opt_passes=5):
    self.update_weights()
    X_b = self.best.turn
    bF = self.best.cost

    for agent in self.agents:
        p = np.tanh(np.abs(agent.cost - bF))
        if random.uniform(0, 1) < p:
            X_A = random.choices(
                self.agents,
                weights=[a.weight for a in self.agents], k=1)[0]
            remaining = [a for a in self.agents if a is not X_A]
            remaining_weights = [a.weight for a in remaining]
            X_B = random.choices(
                remaining, weights=remaining_weights, k=1)[0]
            new_turn = ox_crossover(X_b, ox_crossover(X_A.turn, X_B.turn))
        else:
            new_turn = random_swap(agent.turn)

        new_cost = self.graph.turn_cost(new_turn)
        if new_cost < agent.cost:
            agent.turn = new_turn
            agent.cost = new_cost

        if t % phase_length >= phase_length * (1 - opt_ratio):
            agent.turn = two_opt(
                agent.turn, self.graph, max_passes=max_opt_passes)
            agent.cost = self.graph.turn_cost(agent.turn)
\end{lstlisting}

\section{Opérateurs}
\begin{lstlisting}
def random_swap(turn):
    tmp = turn[:]
    a, b = random.randint(0, len(turn)-1), random.randint(0, len(turn)-1)
    tmp[a], tmp[b] = tmp[b], tmp[a]
    return tmp

def ox_crossover(parent1, parent2):
    n = len(parent1)
    a, b = sorted([random.randint(0, n-1), random.randint(0, n-1)])
    enfant = [None] * n
    enfant[a:b] = parent1[a:b]
    placed = set(parent1[a:b])
    remaining = [v for v in parent2 if v not in placed]
    j = 0
    for i in range(n):
        if enfant[i] is None:
            enfant[i] = remaining[j]
            j += 1
    return enfant

def two_opt(turn, graph, max_passes=5):
    n = len(turn)
    for _ in range(max_passes):
        improved = False
        for i in range(0, n-2):
            for j in range(i+2, n-1):
                gain = (graph.distance(turn[i], turn[i+1])
                      + graph.distance(turn[j], turn[j+1])
                      - graph.distance(turn[i], turn[j])
                      - graph.distance(turn[i+1], turn[j+1]))
                if gain > 0:
                    turn[i+1:j+1] = turn[i+1:j+1][::-1]
                    improved = True
        if not improved:
            break
    return turn
\end{lstlisting}

\section{Boucle d'exécution avec arrêt anticipé}
\begin{lstlisting}
def run(self, T, phase_length=15, opt_ratio=0.25,
        max_opt_passes=5, on_step=None):
    opt_phase_start = int(phase_length * (1 - opt_ratio))
    best_before_opt = None
    best_after_opt = None
    opt_was_stale = False
    history = []

    for t in range(T):
        cycle_pos = t % phase_length
        is_opt_phase = cycle_pos >= opt_phase_start
        phase = "2-opt" if is_opt_phase else "exploration"

        self.step(t, phase_length=phase_length,
                  opt_ratio=opt_ratio, max_opt_passes=max_opt_passes)
        best_cost = self.best.cost
        history.append(best_cost)

        if on_step:
            on_step(t, self, history, phase, stopped=False)

        # Entrée en 2-opt : snapshot du meilleur coût
        if cycle_pos == opt_phase_start:
            best_before_opt = best_cost

        # Pendant le 2-opt : vérifier la stagnation
        if is_opt_phase:
            opt_was_stale = (best_before_opt is not None
                             and best_cost >= best_before_opt)
            best_after_opt = best_cost

        # Exploration suivante : si 2-opt était stale
        # et aucune amélioration -> arrêt
        if not is_opt_phase and opt_was_stale:
            if (best_after_opt is not None
                    and best_cost >= best_after_opt):
                if on_step:
                    on_step(t, self, history, phase, stopped=True)
                break

    return self.best.cost, len(history), history, self.best.turn
\end{lstlisting}

% -----------------------------------------------
\chapter{Benchmark TSP}

\section{Infrastructure de benchmark}
L'implémentation fournit deux outils de benchmark :
\begin{itemize}
	\item \texttt{bench\_single.py} : évaluation sur un ou plusieurs fichiers
	      CSV (fusionnés le cas échéant), avec visualisation en temps réel
	      optionnelle et paramétrage via la ligne de commande ;
	\item \texttt{benchmark.py} : campagne comparative sur plusieurs instances
	      de tailles variées (20, 50, 100 villes issues de TSPLIB, plus une
	      instance \textit{world}).
\end{itemize}

Les deux outils exécutent l'algorithme avec plusieurs tailles de population
($10$, $20$, $40$, $80$ agents) et produisent des graphiques comparatifs
(tournées, courbes de convergence, coût \textit{vs.} temps).

\section{Formats de données acceptés}
\begin{itemize}
	\item \textbf{Format TSP dataset} : colonne \texttt{city\_coordinates}
	      contenant une liste de tuples, et colonne \texttt{total\_distance}
	      pour la solution optimale connue ;
	\item \textbf{Format coordonnées} : colonnes \texttt{x}/\texttt{XCOORD.}
	      et \texttt{y}/\texttt{YCOORD.} (instances Solomon, \textit{world}).
\end{itemize}

\section{Métriques}

\begin{table}[h]
	\centering
	\begin{tabular}{@{}ll@{}}
		\toprule
		\textbf{Métrique} & \textbf{Définition} \\
		\midrule
		\% réussite        & Proportion de runs où le gap avec l'optimal est $\leq 1\,\%$ \\
		\% non-détection   & L'algorithme estime avoir convergé mais le gap réel est $> 1\,\%$ \\
		\% fausse détection & L'algorithme s'arrête en estimant échouer alors que le gap est $\leq 1\,\%$ \\
		Gap moyen          & Écart relatif moyen par rapport à l'optimal sur $N$ exécutions \\
		Écart-type du gap  & Mesure de la robustesse sur $N$ exécutions \\
		Itérations         & Nombre d'itérations avant convergence (ou arrêt anticipé) \\
		Temps              & Durée totale d'exécution (en secondes) \\
		\bottomrule
	\end{tabular}
	\caption{Définition des métriques de benchmark}
\end{table}

\section{Résultats}

\begin{table}[h]
	\centering
	\begin{tabular}{@{}rcccc@{}}
		\toprule
		\textbf{Nœuds} & \textbf{\% réussite} & \textbf{\% non-détection}
					& \textbf{\% fausse détection} & \textbf{Temps (s)} \\
		\midrule
		20         & -- & -- & -- & -- \\
		50         & -- & -- & -- & -- \\
		100        & -- & -- & -- & -- \\
		200        & -- & -- & -- & -- \\
		500        & -- & -- & -- & -- \\
		1\,000     & -- & -- & -- & -- \\
		\bottomrule
	\end{tabular}
	\caption{Benchmark SMA sur instances TSP (à compléter)}
\end{table}

\end{document}
